\documentclass[12pt]{article}
\usepackage[utf8]{inputenc}
\usepackage[T1]{fontenc}
\usepackage{hyperref}
\title{Yucuna Language Resources for K-12 Curriculum Development}
\author{Compiled by Claude Sonnet 5.5 (claude-sonnet-5-5) for Andr\'as Kornai}
\date{}
\begin{document}
\maketitle

\section{Introduction}
Yucuna (also spelled Yukuna or Jukuna; ISO 639-3 ycn; Glottocode yucu1253) is an Arawakan (Maipurean) language of Colombia, with about 1,000 speakers. Speakers live mainly along the Mirit\'i-Paran\'a River. The search found academic descriptions and a documentation corpus. It found no school materials.

\section{Primary Educational Resources}
No Yucuna primers or textbooks were found in the search. A literacy handbook is a stated goal of the ELDP documentation project below, but no handbook was seen.

\section{Linguistic and Reference Materials}
\begin{itemize}
\item Magdalena Lemus Serrano, doctoral dissertation on pervasive nominalization in Yukuna, with a grammar sketch: \url{https://inalco.academia.edu/MagdalenaLemusSerrano} and \url{https://www.academia.edu/44761451}
\item SAPhon phonological inventory for Yucuna, University of California, Berkeley: \url{https://linguistics.berkeley.edu/saphon/en/inv/Yucuna.html}
\item Yucuna alphabet and grammar exhibit: \url{https://iris.siue.edu/eng-318/exhibits/show/yucunalanguage/alphabet-grammer}
\item Wikipedia, Yucuna language: \url{https://en.wikipedia.com/wiki/Yukuna}
\end{itemize}
No standalone published dictionary was found.

\section{Religious and Community Resources}
None were found in the search.

\section{Digital and Audio Resources}
\begin{itemize}
\item ``A preliminary linguistic survey of Yucuna, an endangered language of Colombia'' (ELDP project page): \url{https://cultureincrisis.org/projects/a-preliminary-linguistic-survey-of-yucuna-an-endangered-language-of-colombia}
\item The same project's archive collection of transcribed recordings: \url{https://eldp.access.preservica.com/dk0406/}
\item Omniglot page on Yucuna (Jukuna): \url{https://www.omniglot.com/writing/yucuna.htm}
\end{itemize}

\section{Curriculum Development Context}
Yucuna is stable and probably spoken by all community members, so a native-language curriculum would build on spoken use. The main available base is academic documentation: the grammar sketch, the phonological inventory and the recorded corpus. Primers, a dictionary and an agreed orthography would still have to be developed with the community.
\end{document}